\clearpage
\phantomsection\label{sec:references}
\begingroup\small
\begin{thebibliography}{99}
\bibitem{ref1} M. Welzl, T. Henderson, G. Fairhurst, and M. P. Tahiliani. Increase of the Congestion Window when the Sender Is Rate-Limited. Internet-Draft, \nolinkurl{draft-ietf-ccwg-ratelimited-increase-11}, 6 September 2026. Sections 3 and 4.5. \url{https://www.ietf.org/archive/id/draft-ietf-ccwg-ratelimited-increase-11.html}
\bibitem{ref2} L. Xu, S. Ha, I. Rhee, V. Goel, and L. Eggert. CUBIC for Fast and Long-Distance Networks. RFC 9438, August 2023. Sections 4.2 and 5.8. \url{https://www.rfc-editor.org/rfc/rfc9438.html}
\bibitem{ref3} J. Iyengar and I. Swett. QUIC Loss Detection and Congestion Control. RFC 9002, May 2021. Section 7.8. \url{https://www.rfc-editor.org/rfc/rfc9002.html}
\bibitem{ref4} J. Iyengar and M. Thomson. QUIC: A UDP-Based Multiplexed and Secure Transport. RFC 9000, May 2021. \url{https://www.rfc-editor.org/rfc/rfc9000.html}
\bibitem{ref5} M. Allman, V. Paxson, and E. Blanton. TCP Congestion Control. RFC 5681, September 2009. \url{https://www.rfc-editor.org/rfc/rfc5681.html}
\bibitem{ref6} G. Fairhurst, A. Sathiaseelan, and R. Secchi. Updating TCP to Support Rate-Limited Traffic. RFC 7661, October 2015. Experimental. \url{https://www.rfc-editor.org/rfc/rfc7661.html}
\bibitem{ref7} P. Balasubramanian, Y. Huang, and M. Olson. HyStart++: Modified Slow Start for TCP. RFC 9406, May 2023. \url{https://www.rfc-editor.org/rfc/rfc9406.html}
\bibitem{ref8} Private Octopus. picoquic: official source repository. The local experiment revisions are identified separately in Appendix A. \url{https://github.com/private-octopus/picoquic}
\bibitem{ref9} Linux man-pages project. network\_namespaces(7): overview of Linux network namespaces. \url{https://man7.org/linux/man-pages/man7/network_namespaces.7.html}
\bibitem{ref10} iproute2 manual. tc-netem(8): Network Emulator. \url{https://man7.org/linux/man-pages/man8/tc-netem.8.html}
\bibitem{ref12} Local project sources and implementation records: picoquic/cubic.c, sender.c, picoquic\_internal.h; picoquictest/congestion\_test.c; PROJECT\_PROGRESS.md, IMPLEMENTATION\_PLAN.md, and results/README.md; topology/setup.sh; experiments/collect\_results.py and validate\_rli.py.
\bibitem{ref13} OSNT Project. Ubuntu network-namespace campaign, 10 November 2026. The local \path{results/} bundle contains source identifiers, binary hashes, measurements, figures, and raw qlogs. The previous campaign is retained separately. \url{https://github.com/yoman12357/OSNT_Project}
\end{thebibliography}
\endgroup
